\section{Introduction}\label{sec:intro}

Let $C:\bits^n\to\bits$ be a Boolean circuit and $D$ a distribution on its inputs.
We want $p=\Pr_{x\sim D}[C(x)=1]$ when $p$ is tiny. Monte Carlo sampling finds
it to relative error $\eps$ with $\Theta(\log(1/\delta)/(p\eps^2))$ evaluations
of $C$: polynomial in $1/p$, so exponential in the \emph{surprise}
$\log_2(1/p)$ of the event. This note asks two questions.
\begin{enumerate}
\item \emph{Without help.} For which circuits can a polynomial-time procedure,
one that may check some property of the circuit and then sample, estimate $p$
to relative error at a cost polynomial in $\log(1/p)$ instead of $1/p$, while
never being wrong (beyond probability $\delta$) on any circuit?
\item \emph{With a prover.} An unbounded prover sends one argument; a
polynomial-time verifier checks it and samples. No argument may make the
verifier accept a false claim with probability above $\delta$, on any circuit.
For which circuits do convincing arguments exist?
\end{enumerate}
These are certified versions of the low-probability estimation problems studied
by the Alignment Research Center \cite{wu2024lowprob,christiano2022presumption},
where the aim is to estimate the probability of a rare behaviour of a model
faster than by sampling, using the model's internals.

\subsection{The picture}

One idea organises almost everything: \textbf{upper bounds need deduction;
lower bounds need proposals; sampling pays for whatever surprise is left}.
\begin{itemize}
\item A lower bound $p\ge L$ can be certified by exhibiting mass: a proposal
distribution concentrated where $C=1$, with computable likelihood ratios.
Clipped importance weights make this sound for any proposal (\cref{prop:proposal}).
\item An upper bound below the sampling scale cannot be certified by
evaluating $C$, however cleverly: a verifier that makes $T$ evaluations cannot
certify $p<1/T$ (up to constants), whatever argument it is given
(\cref{thm:blackbox,thm:surprise}). It must read the circuit and \emph{deduce}
that $C(x)=1$ confines $x$ to a region of small, exactly computable measure.
\item Combining the two, an \emph{envelope}, a tractable function $\varphi$ with
$\varphi\ge\one_C$ whose mass $Z=\E_D\varphi$ is computable and from which one can
sample, yields an importance sampler with bounded weights, hence two-sided
intervals, at cost $O(K\eps^{-2}\log(1/\delta))$ where $K=Z/p$
(\cref{thm:envelope}). The envelope explains $\log_2(1/Z)$ bits of the
surprise; the remaining $\log_2K$ bits are paid for by sampling, and this
residual cost is unavoidable for anything that treats the unexplained part as
a black box (\cref{thm:surprise}).
\end{itemize}
Karp--Luby--Madras estimation of DNF probabilities \cite{klm1989}, exact
counting for decision-DNNFs and bounded treewidth, exponential tilting for
large deviations \cite{siegmund1976} and the Chernoff bound are all instances.

\subsection{Results}

\paragraph{Limits (\cref{sec:limits}).} Beating sampling on all circuits is
impossible: it would put $\NP$ in $\BPP$, and with exponentially hard
pseudorandom generators no procedure beats sampling by more than a $p^{-o(1)}$
factor even on explicit small circuits (\cref{prop:worstcase}). With a prover, no
argument system is complete for upper bounds unless $\coNP\subseteq\MA$
(\cref{thm:ma}), and black-box verification of upper bounds is impossible
(\cref{thm:blackbox}). Interaction would trivialise the problem (sum-check
gives exact $p$ for every circuit), so the one-message model is the
interesting one.

\paragraph{Without a prover (\cref{sec:part1}).} The procedure DFSB
(Decompose--Force--Sample--Branch) recursively splits the event by exact rules
(independence of variable-disjoint parts, literals forced by propagation,
Shannon expansion) and samples only at leaves that are dense or have tight
envelopes. It is sound on every circuit, and complete, with cost
$O(m\,K\,d^2\eps^{-2}\log(m/\delta))$ evaluations, on every circuit whose ideal
$K$-explanation tree has $m$ nodes (\cref{thm:dfsb}). This covers products and
read-once compositions of arbitrary dense sub-circuits, backdoor-style triggers,
circuits with small cutsets (exactly, by cutset conditioning), DNFs and DNNFs, and,
through sound local rewriting, disguised versions of all of these. Specialised
modules add bounded treewidth (exact), thresholds of independent blocks (exponential
tilting with $K=O(\sqrt k)$, \cref{prop:tilting}), unate CNFs with sparse overlaps (a
deterministic Harris--Janson sandwich, \cref{thm:janson}, which certifies the
``presumption of independence''), and the known FPRAS classes
\cite{meel2026dnnf,moitra2019,gopalan2011fptas}.

\paragraph{With a prover (\cref{sec:part2}).} Proof-carrying importance sampling:
the argument is a DAG of deductive steps checked by any Cook--Reckhow proof
system (implications, exclusivity, independence) and sampling steps with
declared confidence budgets. It is sound on every circuit and every argument
(\cref{thm:pcis-sound}), and complete on every circuit that has a short proof of
being confined to a tractable region in which it is not rare
(\cref{thm:certenv}). Lower bounds need no proof at all. The verifier's cost is
$\poly(|\pi|)+2^{\sigma_u}$, where $\sigma_u$ is the surprise left unexplained by
the argument, and both terms are necessary in the black-box sense: each bit of
argument buys at most one bit of surprise (\cref{thm:surprise}). The prover helps
strictly: for lower bounds unless $\NP\subseteq\BPP$; for upper bounds exactly to
the extent that the verifier's proof system is not weakly automatable
(\cref{cor:automatability}); two-sidedly under the hardness of factoring. Such
separations cannot be based on $\NP$-hardness alone unless $\coNP\subseteq\AM$.

\paragraph{Experiments (\cref{sec:experiments}).} A Python implementation
(\texttt{rareprob}) estimates, e.g., $p\approx4\cdot10^{-44}$ for a threshold of 400
independent random blocks to $5\%$ with $5\cdot10^5$ samples, $p\approx7\cdot10^{-25}$
for an 80-bit trigger with $2\cdot10^4$ samples, and certifies a
$2^{-20}$-rare SPN-hidden trigger with a SAT-solver proof where all syntactic methods
fail; in 2{,}450 adversarial and random trials no interval missed $p$ more often
than allowed.

\subsection{Related work}

\emph{Approximate counting.} Karp--Luby--Madras \cite{klm1989} and the stopping
rule of \cite{dklr2000} are the prototypes; for self-reducible classes, approximate
counting and approximate sampling are equivalent \cite{jvv1986}, and Karger's
unreliability FPRAS \cite{karger1999} is an early instance of a combinatorial
theorem (few near-minimum cuts) serving as the envelope. FPRASs exist for DNNF
\cite{meel2026dnnf} (and before it nFBDD \cite{meel2025nfbdd} and NFA
\cite{arenas2021nfa}), for CNFs in the local-lemma regime
\cite{moitra2019,fgyz2021,fhy2021,jpv2021,hwy2022,wangyin2024}, and for knapsack
\cite{gopalan2011fptas,gkm2010}; d-DNNF, OBDD, SDD and bounded treewidth admit exact
counting \cite{darwiche2002map,samer2010}. Hashing-based counters
\cite{chakraborty2013approxmc} use an NP oracle; partial compilation with
sampling \cite{lai2023partialkc} and AND/OR importance sampling
\cite{gogate2008andor} combine decomposition with sampling for variance reduction,
without relative-error guarantees for rare events. Certified exact model counting
\cite{bryant2023cpog} checks compilation by proofs.
\emph{Rare-event simulation} \cite{asmussenglynn2007,bucklew2004} studies
importance sampling with bounded relative error; the soundness requirement
here (valid intervals on every input) is the computer-science twist.
\emph{Interactive proofs} certify exact counts \cite{lfkn1992} and lower bounds
\cite{gs1986}; the Merlin--Arthur setting connects to proof complexity
\cite{cookreckhow1979,kp1998,bpr2000,pudlak2003pairs} and to the class $\SBP$
\cite{bgm2006}.
\emph{Heuristic arguments.} ARC's programme
\cite{christiano2022presumption,hilton2024surprise,wu2024lowprob} seeks estimators
that use a model's internals to beat sampling, judged by ``surprise accounting'';
this note isolates the part of that programme that can be made sound, and shows
where soundness forces deduction. ARC's ``matching sampling principle''
\cite{neyman2025competing} asks for estimators that do at least as well as sampling
(on average over a distribution of circuits, without soundness), and its
``computational no-coincidence principle'' \cite{neyman2025nocoincidence}
conjectures that structure in random reversible circuits always has short
explanations checkable by a verifier; our Part~2 is a worst-case, statistically
sound relative of the latter, for probabilities instead of properties.
